% Generated by build/make_tables.py from regime_table.csv. Do not edit.

\begin{table}[tbp]
\centering
\footnotesize
\setlength{\tabcolsep}{4.5pt}
\caption{The fixed book across market regimes}
\label{tab:regimes}
\begin{tabular}{@{}lrrrrrrr@{}}
\toprule
 & & Mean & Diversif. & Realised & Exception & Equity vol & 10y vol \\
Regime & Days & VaR & (\%) & P\&L s.d. & rate (\%) & (ann., \%) & (bp/day) \\
\midrule
Calm, 2017--2019 & 687 & 54.8 & 44.8 & 28.7 & 3.5 & 12.7 & 4.0 \\
COVID shock, 2020 & 82 & 56.8 & 45.4 & 67.9 & 11.0 & 57.2 & 8.8 \\
Recovery, 2021 & 231 & 52.9 & 45.8 & 31.0 & 4.8 & 13.7 & 4.5 \\
Rate shock, 2022 & 195 & 64.9 & 42.3 & 94.9 & 32.3 & 25.4 & 8.4 \\
After the shock, 2024--2025 & 452 & 99.7 & 32.1 & 49.2 & 3.3 & 16.7 & 5.8 \\
\bottomrule
\end{tabular}
\notetext{The book is sized once on 2019 and never changed, so every difference across rows is the market rather than the book. Mean VaR and realised profit-and-loss standard deviation are in millions of US dollars at the book's normalising scale and carry no information in levels; their ratio within a row does. The last row is where the downward drift in the measured diversification effect shows up, with VaR at its sample peak and volatility back to ordinary levels.}
\end{table}
